% v1-H: Correlation-Based Change-Point Detector — Mathematical formulation
% Author: Mario de los Santos
% Date: 2026-05-26
%
% Build:   pdflatex 2026-05-26_v1H_correlation_detector_math.tex
% Requires: amsmath, amssymb, mathtools, algorithm, algpseudocode, booktabs, hyperref, microtype

\documentclass[11pt,a4paper]{article}

\usepackage[utf8]{inputenc}
\usepackage[T1]{fontenc}
\usepackage{lmodern}
\usepackage[margin=1in]{geometry}
\usepackage{amsmath,amssymb,amsthm,mathtools}
\usepackage{algorithm}
\usepackage{algpseudocode}
\usepackage{booktabs}
\usepackage{microtype}
\usepackage[hidelinks]{hyperref}
\usepackage{xcolor}

\newcommand{\Var}{\operatorname{Var}}
\newcommand{\Cov}{\operatorname{Cov}}
\newcommand{\corr}{\operatorname{corr}}
\newcommand{\MAD}{\operatorname{MAD}}
\newcommand{\med}{\operatorname{median}}
\newcommand{\vectri}{\operatorname{vec}_{\triangle}}
\newcommand{\R}{\mathbb{R}}
\newcommand{\E}{\mathbb{E}}
\newcommand{\Real}{\mathbb{R}}

\title{\textbf{v1-H: A Correlation-Based Change-Point Detector\\for Non-Stationary Causal Discovery}\\[4pt]
\large Mathematical formulation}
\author{Mario de los Santos}
\date{2026-05-26}

\begin{document}
\maketitle

\begin{abstract}
We present \textsf{v1-H}, a sliding-window change-point detector that operates on the
inter-channel Pearson correlation structure of a multivariate time series.
The detector targets a specific class of structural change---causal-graph
rewiring---for which the published wavelet multi-moment baseline
(\textsf{v1-G}) is shown empirically to collapse at long regime lengths.
Building on the second-order break-detection framework of
Aue, Hörmann, Horváth \& Reimherr (2009), we adapt the
statistic from covariance to Pearson correlation for scale invariance,
replace asymptotic CUSUM testing with localized peak detection on a
per-time-step Frobenius distance, and calibrate thresholds robustly via
the median absolute deviation. This note documents the mathematical
formulation: notation, the change signal, calibration, peak extraction,
complexity, and an algebraic argument for why a correlation-shift statistic
detects exactly the class of changes that a per-channel-moment statistic
misses.
\end{abstract}

\section{Notation and setup}

Let $X \in \R^{T \times N}$ denote a discrete-time multivariate process with
$T$ samples and $N \ge 2$ channels. We model $X$ as a piecewise-stationary
process: there exist unknown change points
$0 < \tau_1 < \tau_2 < \cdots < \tau_K < T$ which partition the time axis
into $K+1$ \emph{regimes}
$R_0 = [0, \tau_1), R_1 = [\tau_1, \tau_2), \ldots, R_K = [\tau_K, T)$.
Within each $R_k$ the joint distribution of $X[t,\cdot]$ is stationary; at
each $\tau_k$ the underlying structural causal model (SCM) parameters change.

In the linear non-Gaussian acyclic SCM framework (LiNGAM,
Shimizu \emph{et al.}\ 2006), the parameters of regime $R_k$ consist of a
directed acyclic graph $G_k$ over $\{V_1, \ldots, V_N\}$ together with
edge weights $B_k \in \R^{N \times N}$, such that
\[
X[t, \cdot]^\top = B_k\, X[t, \cdot]^\top + \varepsilon^{(k)}, \qquad t \in R_k,
\]
with $\varepsilon^{(k)}$ a vector of mutually independent non-Gaussian
exogenous noises. A structural change at $\tau_k$ corresponds to
$B_{k-1} \ne B_k$, typically the addition or removal of one or more directed
edges $V_i \to V_j$.

\paragraph{Detection task.}
Given $X$ and a baseline index set $\mathcal{B} \subset \{0, 1, \ldots, T-1\}$
known to lie within $R_0$, output a detected change-point set
$\widehat{\mathcal{T}} = \{\hat\tau_1, \ldots, \hat\tau_{\hat K}\}$.

\section{The change signal}

\subsection{Pairwise correlation matrices}

For a window size $W$ and an interior time $t \in [W,\, T - W)$, define the
\emph{past} and \emph{future} windows
\begin{equation}
X_{\text{past}}(t) = X[t - W : t,\ \cdot]
\qquad
X_{\text{future}}(t) = X[t : t + W,\ \cdot].
\end{equation}

The associated Pearson correlation matrices
$C_{\text{past}}(t), C_{\text{future}}(t) \in \R^{N \times N}$
have entries
\begin{equation}
C_{\text{past}}(t)_{ij}
= \frac{\sum_{s = t - W}^{t - 1}\!\big(X[s,i] - \bar X_i^{\text{past}}\big)\big(X[s,j] - \bar X_j^{\text{past}}\big)}
{\,\hat\sigma_i^{\text{past}}\, \hat\sigma_j^{\text{past}}\,},
\end{equation}
where $\bar X_i^{\text{past}}$ and $\hat\sigma_i^{\text{past}}$ are the sample
mean and standard deviation of channel $i$ in $X_{\text{past}}(t)$; the
definition of $C_{\text{future}}(t)$ is analogous. Both matrices are
symmetric with unit diagonal, so the structural information they carry
lives in the $M = \binom{N}{2}$ strictly-upper-triangular entries. We
vectorize:
\begin{equation}
c_{\text{past}}(t) = \vectri\!\big(C_{\text{past}}(t)\big) \in \R^{M},
\qquad
c_{\text{future}}(t) = \vectri\!\big(C_{\text{future}}(t)\big) \in \R^{M}.
\end{equation}

\subsection{Definition of the change signal}

The change signal at time $t$ is the Euclidean distance between past and
future correlation fingerprints:
\begin{equation}
\boxed{\;
S(t) \;=\; \big\| c_{\text{future}}(t) - c_{\text{past}}(t) \big\|_2
\;=\; \sqrt{\sum_{1 \le i < j \le N}\!\big(C_{\text{future}}(t)_{ij} - C_{\text{past}}(t)_{ij}\big)^2}.
\;}
\label{eq:change-signal}
\end{equation}

Equivalently, in terms of the Frobenius norm of the matrix difference
(noting that the off-diagonal entries appear twice and the diagonal is
identically zero since $C_{ii} \equiv 1$):
\begin{equation}
S(t) \;=\; \frac{1}{\sqrt{2}}\, \big\| C_{\text{future}}(t) - C_{\text{past}}(t) \big\|_F.
\end{equation}

\subsection{Properties of $S(t)$}

\paragraph{P1.\ Vanishing under stationarity.}
If the joint distribution of $X$ is stationary on $[t - W, t + W]$, then
$c_{\text{past}}(t)$ and $c_{\text{future}}(t)$ are i.i.d.\ unbiased estimates
of the population correlation vector $c^\star$. By the multivariate
central limit theorem applied to the sample correlations,
\begin{equation}
\sqrt{W}\,\big(c_{\text{past}}(t) - c^\star\big)
\xrightarrow{\,d\,} \mathcal{N}(0,\ \Sigma_c),
\end{equation}
with $\Sigma_c$ depending on the population fourth-order moments. Hence
$\big\| c_{\text{future}}(t) - c_{\text{past}}(t) \big\|_2$ is
$\mathcal{O}_p(W^{-1/2})$ under the null hypothesis of local stationarity.

\paragraph{P2.\ Localized peak at change points.}
Suppose $\tau \in (W, T - W)$ is a change point and $W$ is small enough that
$X_{\text{past}}(\tau)$ lies entirely within $R_{k-1}$ and
$X_{\text{future}}(\tau)$ entirely within $R_k$. Then
\begin{equation}
\E\!\big[S(\tau)\big] \;\to\; \big\| c^{(k)} - c^{(k-1)} \big\|_2
\quad \text{as } W \to \infty,
\end{equation}
the deterministic Euclidean distance between the regime-$k-1$ and regime-$k$
population correlation vectors. For $|t - \tau| > W$, both windows lie in
the same regime and $S(t) = \mathcal{O}_p(W^{-1/2})$. Therefore $S$ exhibits
a peak with finite expected height (independent of $W$, for $W$ large) at
each $\tau_k$, against a background of $\mathcal{O}(W^{-1/2})$ noise.

\paragraph{P3.\ Affine scale invariance.}
Under the channelwise affine transformation $X_i \mapsto a_i X_i + b_i$
with $a_i \ne 0$, both $C_{\text{past}}(t)$ and $C_{\text{future}}(t)$ are
invariant, hence so is $S(t)$. This property is critical for the LiNGAM
setting where the per-regime noise variance varies as a design choice;
covariance-based statistics would conflate this with structural change.

\section{Why correlation detects what per-channel moments miss}

We give an algebraic argument for the central claim of v1-H: that
\emph{causal-graph rewiring} produces a first-order shift in the
inter-channel correlation but only a second-order shift in per-channel
moments. Consider a linear-Gaussian SCM with all channels normalized to
unit variance, and suppose a single edge $V_i \to V_j$ of weight $\beta$
is added at change point $\tau$. The reduced form for $V_j$ at $t < \tau$ is
\begin{equation}
V_j = \sum_{k \in \mathrm{pa}_{j,k-1}} \beta_{kj}\, V_k + \varepsilon_j,
\end{equation}
and at $t > \tau$ gains an additional term:
\begin{equation}
V_j = \beta\, V_i + \sum_{k \in \mathrm{pa}_{j,k-1}} \beta_{kj}\, V_k + \varepsilon_j.
\end{equation}

\paragraph{Per-channel mean.}
\begin{equation}
\E[V_j \mid t > \tau] - \E[V_j \mid t < \tau] \;=\; \beta \cdot \E[V_i].
\end{equation}
For zero-mean centred channels, $\E[V_i] = 0$, so the marginal-mean shift is
\emph{exactly zero} regardless of $\beta$.

\paragraph{Per-channel variance.}
\begin{align}
\Delta \Var(V_j)
&= \Var(V_j \mid t > \tau) - \Var(V_j \mid t < \tau) \nonumber \\
&= \beta^2 \Var(V_i) + 2\beta \sum_{k \in \mathrm{pa}_{j,k-1}} \beta_{kj}\,\Cov(V_k, V_i).
\end{align}
The leading term is $\mathcal{O}(\beta^2)$, modulated by the existing
covariance structure.

\paragraph{Pairwise correlation.}
Under the unit-variance normalization (which is achieved post-hoc by the
correlation operation itself):
\begin{equation}
\corr(V_i, V_j \mid t > \tau) - \corr(V_i, V_j \mid t < \tau)
\;\approx\; \beta \cdot \sqrt{\frac{\Var(V_i)}{\Var(V_j)}}\,.
\end{equation}
The shift is $\mathcal{O}(\beta)$, with a constant of order $1$ for
similarly-scaled channels.

\paragraph{Implication.}
A detector that integrates per-channel moments---such as the multi-moment
wavelet detector (v1-G), which sums squared wavelet coefficients of mean,
variance, skewness, and kurtosis time series---has sensitivity proportional
to the $\mathcal{O}(\beta^2)$ variance shift and smaller-order shifts in
the higher moments. The correlation-shift statistic in
Eq.~\eqref{eq:change-signal} has sensitivity proportional to the
$\mathcal{O}(\beta)$ correlation shift. For typical edge weights
$\beta \in [0.5, 1.5]$ in the benchmark, this is an order-of-magnitude
advantage at the change point.

\section{MAD-based threshold calibration}

Let $\mathcal{B} \subset \{0, 1, \ldots, T-1\}$ be a baseline index set
restricted to within-regime indices in $R_0$. Define
\begin{align}
\mu_{\text{med}} &= \med\!\big(\{S(t) : t \in \mathcal{B}\}\big), \\
\MAD\,(\mathcal{B}) &= \med\!\big(\{|S(t) - \mu_{\text{med}}| : t \in \mathcal{B}\}\big), \\
\hat\sigma_S &= 1.4826 \cdot \MAD\,(\mathcal{B}).
\end{align}

The scale factor $1.4826 = 1/\Phi^{-1}(0.75)$, where $\Phi$ is the standard
normal CDF, ensures that $\hat\sigma_S$ is an asymptotically unbiased
estimator of $\mathrm{SD}(S)$ when $S$ is Gaussian on $\mathcal{B}$.

The peak-detection threshold is
\begin{equation}
\boxed{\;
\theta \;=\; \max\!\big(\,\mu_{\text{med}} + k \cdot \hat\sigma_S,\;\; \tau_{\min}\,\big),
\;}
\end{equation}
with strictness hyperparameter $k > 0$ and floor $\tau_{\min} \ge 0$.
v1-H uses $k = 6.0$ and $\tau_{\min} = 0.5$.

\paragraph{Why MAD rather than mean$\,+\,$standard deviation?}
The mean$\,+\,$SD estimator has a breakdown point of $0$: a single outlier
in $\mathcal{B}$ (for example a missed change point inadvertently included in
the baseline) inflates $\hat\sigma_S$ arbitrarily. MAD has breakdown point
$0.5$. This robustness matters because we cannot test the
baseline-stationarity assumption ahead of time.

\paragraph{Why MAD rather than the asymptotic distribution of Aue et al.?}
The CUSUM-based test in Aue \emph{et al.}\ (2009) derives its threshold
from the asymptotic null distribution of the cumulative-deviation
statistic, under regularity assumptions (mixing rates, finite higher-order
moments). These assumptions are non-trivial to verify in practice. MAD
calibration is fully empirical: it adapts to whatever noise level the
particular series exhibits, without distributional commitments.

\section{Peak detection and refractory enforcement}

Given $\theta$ and the change signal $S$, candidate change points are
strict local maxima of $S$ exceeding $\theta$:
\begin{equation}
\widehat{\mathcal{T}}_0
= \big\{\,t \in [W, T-W) :\ S(t-1) < S(t),\ S(t+1) < S(t),\ S(t) \ge \theta\,\big\}.
\end{equation}

A refractory constraint $\rho > 0$ enforces a minimum temporal separation
between accepted change points, preventing a single physical CP from being
detected as several adjacent peaks. We greedily retain the strongest peak
within each window of width $\rho$:
\begin{equation}
\widehat{\mathcal{T}}
= \mathrm{NMS}_{\rho}\!\big(\widehat{\mathcal{T}}_0;\, S\big),
\end{equation}
where $\mathrm{NMS}_{\rho}$ denotes the standard non-maximum-suppression
operation: iteratively select the highest-$S$ peak not yet selected, then
discard all candidates within $[\hat\tau - \rho, \hat\tau + \rho]$.
This is the canonical implementation provided by
\texttt{scipy.signal.find\_peaks(S, height=$\theta$, distance=$\rho$)}.

\section{Window-size and refractory scaling}

Rather than fix $W$ and $\rho$ as global constants, we scale both with the
minimum expected regime length $L_{\min}$:
\begin{align}
W   &= \max\!\big(50,\ \lfloor L_{\min}/5 \rfloor\big), \\
\rho &= \max\!\big(150,\ \lfloor L_{\min}/4 \rfloor\big).
\end{align}

This produces three operational regimes in the benchmark:
\begin{center}
\begin{tabular}{rrrrr}
\toprule
$L_{\min}$ & $W$ & $\rho$ & samples in $S$ & max $\hat K$ \\
\midrule
500   & 100 & 150  & $\sim T - 200$  & $\le T / 150$ \\
1500  & 300 & 375  & $\sim T - 600$  & $\le T / 375$ \\
2500  & 500 & 625  & $\sim T - 1000$ & $\le T / 625$ \\
\bottomrule
\end{tabular}
\end{center}

The window must be large enough that the empirical correlation estimates
have low variance (Property P2) but small enough to keep the regime-pair
assumption clean (the past window must lie in one regime, the future
window in the next).

\section{Algorithm}

\begin{algorithm}[H]
\caption{v1-H: correlation-based change-point detection}
\label{alg:v1h}
\begin{algorithmic}[1]
\Require Multivariate series $X \in \R^{T \times N}$; baseline indices
         $\mathcal{B}$; minimum regime length $L_{\min}$; strictness $k$;
         floor $\tau_{\min}$
\Ensure  Detected change-point set $\widehat{\mathcal{T}}$
\Statex
\State $W \gets \max(50,\ \lfloor L_{\min}/5 \rfloor)$
\State $\rho \gets \max(150,\ \lfloor L_{\min}/4 \rfloor)$
\State $S \gets \mathbf{0}_T$
\For{$t = W$ \textbf{to} $T - W - 1$}
   \State $C_p \gets \mathrm{corrcoef}\big(X[t-W:t,\cdot]\big)$
   \State $C_f \gets \mathrm{corrcoef}\big(X[t:t+W,\cdot]\big)$
   \State $\Delta c \gets \vectri(C_f) - \vectri(C_p)$ \Comment{$M = \binom{N}{2}$ entries}
   \State $S[t] \gets \|\Delta c\|_2$
\EndFor
\Statex
\State $\mu \gets \med\!\big(\{S[t] : t \in \mathcal{B}\}\big)$
\State $\hat\sigma \gets 1.4826 \cdot \med\!\big(\{|S[t] - \mu| : t \in \mathcal{B}\}\big)$
\State $\theta \gets \max\!\big(\mu + k\,\hat\sigma,\ \tau_{\min}\big)$
\Statex
\State $\widehat{\mathcal{T}} \gets \mathrm{find\_peaks}(S,\ \text{height} = \theta,\ \text{distance} = \rho)$
\State Discard $\hat\tau \in [0, W) \cup [T - W, T)$
\State \Return $\widehat{\mathcal{T}}$
\end{algorithmic}
\end{algorithm}

\section{Computational complexity}

The dominant cost is the per-time-step correlation-matrix computation in
the main loop (lines 4--8 of Algorithm~\ref{alg:v1h}):
\begin{equation}
\mathcal{O}\big(T \cdot W \cdot N^2\big).
\end{equation}
Each of the $T - 2W$ interior time steps reduces a $W \times N$ data block
to an $N \times N$ correlation matrix, which costs $\mathcal{O}(W N^2)$ via
the standard outer-product formulation. For our benchmark
($T \le 1.5 \times 10^4$, $W \le 500$, $N \le 12$), a single scenario runs
in under one second on a single CPU core.

For very long series, an $\mathcal{O}(T N^2)$ implementation is achievable
by maintaining running sums of $X_i$, $X_j$, and $X_i X_j$ across the
sliding window and updating them incrementally as $t$ advances. We have
not implemented this optimization because the benchmark does not require it.

\section{Differences from Aue \emph{et al.}\ (2009)}

\begin{table}[ht]
\centering
\small
\begin{tabular}{p{0.32\linewidth} p{0.30\linewidth} p{0.30\linewidth}}
\toprule
\textbf{Aspect} & \textbf{Aue \emph{et al.}\ (2009)} & \textbf{v1-H} \\
\midrule
Second-order object &
  Empirical covariance &
  Pearson correlation \\
Scale invariance &
  No (depends on channel scaling) &
  Yes (Property P3) \\
Test statistic &
  CUSUM of partial sums &
  Pointwise Frobenius distance (Eq.~\ref{eq:change-signal}) \\
Detection scheme &
  Asymptotic threshold $+$ binary segmentation &
  MAD-calibrated peak detection (single pass) \\
Multiple change points &
  Recursive splits &
  Native via \texttt{find\_peaks} \\
Threshold derivation &
  Asymptotic null distribution &
  Empirical via MAD on baseline \\
Distribution assumption &
  Mixing, finite higher moments &
  Robust to baseline outliers; assumption-free \\
Window size &
  $T$-asymptotic, not specified for finite $T$ &
  Empirically scaled with $L_{\min}$ \\
\bottomrule
\end{tabular}
\end{table}

The two methods share the conceptual core: second-order structural shifts
in a multivariate time series are diagnostic of regime change. v1-H trades
the asymptotic theory of Aue \emph{et al.}\ for empirical robustness, single-pass
operation, and scale invariance, all of which are properties we require
for the detector to operate inside a downstream Bayesian causal-discovery
pipeline.

\section{Recommended hyperparameters}

\begin{center}
\small
\begin{tabular}{lll}
\toprule
\textbf{Parameter} & \textbf{Value} & \textbf{Rationale} \\
\midrule
\texttt{threshold\_mad\_k} ($k$) & 6.0 & Strict threshold; prioritizes precision over recall \\
\texttt{min\_threshold} ($\tau_{\min}$) & 0.5 & Absolute floor; rejects near-zero baseline drift \\
\texttt{window} ($W$) & $\max(50,\ L_{\min}/5)$ & Balances variance vs.\ localization \\
\texttt{refractory\_period} ($\rho$) & $\max(150,\ L_{\min}/4)$ & Matches downstream pipeline's expected regime length \\
\texttt{edge\_margin} & $W$ & Excludes boundary-incomplete windows \\
\bottomrule
\end{tabular}
\end{center}

\section*{Acknowledgements}

The author thanks the supervising committee for guidance on framing
empirical comparisons against the published CPD literature. All numerical
results in companion documents were obtained on a controlled benchmark
described in \texttt{generate\_datasets.py}; the implementation is
in \texttt{NSD\_Wavelets/src/detectors/detectors\_correlation.py}.

\begin{thebibliography}{9}

\bibitem{aue2009}
A.~Aue, S.~H\"ormann, L.~Horv\'ath, and M.~Reimherr.
\newblock Break detection in the covariance structure of multivariate time series.
\newblock \emph{Journal of Multivariate Analysis}, 2009.

\bibitem{cribben2012}
I.~Cribben, R.~Haraldsdottir, L.~Atlas, T.~Wager, and M.~Lindquist.
\newblock Dynamic connectivity regression: determining state-related changes
  in brain connectivity.
\newblock \emph{NeuroImage}, 2012.

\bibitem{truong2020}
C.~Truong, L.~Oudre, and N.~Vayatis.
\newblock Selective review of offline change point detection methods.
\newblock \emph{Signal Processing}, 2020.

\bibitem{harchaoui2009}
Z.~Harchaoui, F.~Bach, and \'E.~Moulines.
\newblock Kernel change-point analysis.
\newblock \emph{Advances in Neural Information Processing Systems}, 2009.

\bibitem{shimizu2006}
S.~Shimizu, P.~O.~Hoyer, A.~Hyv\"arinen, and A.~Kerminen.
\newblock A linear non-Gaussian acyclic model for causal discovery.
\newblock \emph{Journal of Machine Learning Research}, 7:2003--2030, 2006.

\end{thebibliography}

\end{document}
